\documentclass[english,handout]{mlutalk}

\title{Stance Classification for Key-Point Analysis}
% \title{%
%   Modern Talking: Key-Point Analysis \\
%   using Modern Natural Language Processing
% }
\subtitle{Natural Language Processing, Summer Semester 2021}
\author{Max Henze \and Hanh Luu \and Jan Heinrich Reimer}
\institute{Martin Luther University Halle-Wittenberg}
\date{\today}
\titlegraphic{\includegraphics[width=3cm]{figures/mlu-halle}}

\addbibresource{../literature/literature.bib}

\usepackage{tikz}
\usepackage{listings}
\usepackage{xspace}
\usepackage{biblatex}

\newcommand{\Bert}{\textsc{Bert}\xspace}
\newcommand{\ArgKP}{\mbox{ArgKP}\xspace}
\newcommand{\ArgQ}{\mbox{IBM-ArgQ-Rank-30kArgs}\xspace}

\begin{document}

\titleframe

\begin{frame}{Stance Classification}
  TODO: motivation and problem description.
\end{frame}

\begin{frame}{Approaches}
  TODO: general framework and basic approaches.
\end{frame}

\begin{frame}{Fake News Challenge Stance Detection~\cite{HanselowskiSSCC2018}}
   Challenge:
  \begin{itemize}
      \item Stance $s$: argee, disgree, discuss, unrelated
      \item Predict a label stance for a document with regard to a headline: $f: (d,h) \rightarrow s$
  \end{itemize}
  Topic of this paper:
  \begin{itemize}
    \item Reproduction and Analysis for the top-performing systems of the challenge: \textit{Tablos, Athene, UCL}
    \item Question: Which features and architectures used help improving performance?
    \item (?) Imbalanced Learning, in-domain and cross-domain (transfer)
  \end{itemize}
\end{frame}

\begin{frame}{Fake News Challenge Stance Detection~\cite{HanselowskiSSCC2018}}
    \begin{table}[]
    \caption{The best approaches}
        \label{tab:comparison}
        \begin{tabular}{|l|l|l|l|}
        \hline
        \multicolumn{2}{|l|}{Talos} &
          Athene &
          UCL \\ \hline
        \multicolumn{2}{|l|}{\begin{tabular}[c]{@{}l@{}}weighted model\\ (CNN un GBD Tree)\end{tabular}} &
          \begin{tabular}[c]{@{}l@{}}MLP \\ (6 h-layers+ Softmax)\end{tabular} &
          MLP (1 h-layer) \\ \hline
        word2vec &
          \begin{tabular}[c]{@{}l@{}}TF\\ TFIDF\\ Sentiment \\ word2vec\end{tabular} &
          \begin{tabular}[c]{@{}l@{}}Unigrams\\ Similarity (Noun, Verbs)\\ Topic Models\end{tabular} &
          \begin{tabular}[c]{@{}l@{}}TF of unigrams\\ TFIDF \\ Similarity\end{tabular} \\ \hline
        \end{tabular}
    \end{table}
\end{frame}

\begin{frame}{Fake News Challenge Stance Detection~\cite{HanselowskiSSCC2018}}
   Which features and models help improving performance?
   ...
\end{frame}

\begin{frame}{Argument Stance Classification Using \Bert~\cite{OllingerDSBS2020}}
  TODO
\end{frame}

\begin{frame}{Cross-topic Argument Mining~\cite{StabMSRG2018}}
  Topic of this paper:
  \begin{itemize}
    \item To search for arguments, relevant to a certain topic 
    \item Given a heterogeneous document collection
  \end{itemize}
  Key Achievements found in this paper:
  \begin{itemize}
    \item Non, supporting and opposing styled annotation scheme
    \item New corpus including 27,520 sentences for 8 controversial topics
    \item Contextual BiLSTM leads to better generalization to unknown topics
    \item Multi-task learning setup further improved this approach 
  \end{itemize}
\end{frame}

\begin{frame}{Cross-topic Argument Mining~\cite{StabMSRG2018}}
  How did they achieve this ?\\
  $\rightarrow$ two major approaches for identifying arguments
  \begin{itemize}
    \item Integrating topic information
    \begin{itemize}
      \item Outer-attention BiLSTM 
      \begin{itemize}
        \item Model learning how to weight input words given the topic combined with BiLSTM
      \end{itemize}
      \item Custom contextual BiLSTM
      \begin{itemize}
        \item Adding topic information as additional input term 
      \end{itemize}
    \end{itemize}
    \item Leveraging additional data (two additional corpora)
    \begin{itemize}
      \item Transfer learning
      \begin{itemize}
        \item train model twice with auxiliary corpus first 
      \end{itemize}
      \item Multi-task learning
      \begin{itemize}
        \item using RNN for separate learning and combination afterwards
      \end{itemize}
    \end{itemize}
  \end{itemize}
\end{frame}

\begin{frame}{Comparison}
  TODO: compare approaches
\end{frame}

\begin{frame}{Stance in Key Point Matching}
  TODO: how can we use stance classification for matching key points?
\end{frame}

\begin{frame}{Conclusion}
  TODO: what to do next?
\end{frame}

\begin{frame}{Future Work}
  TODO: do we need future work?
  \thankyou
\end{frame}

\appendix
\section{\appendixname}

\bibliographyframe

\end{document}
